\section{练习}

\begin{question}
\begin{drill}[空气与玻璃界面的两次入射]\label{d:11}
空气中的光以 \(45^\circ\) 入射到折射率 \(n=\sqrt2\) 的玻璃，求折射角。另有一束光从同种玻璃射向空气，入射角为 \(50^\circ\)，判断能否有传播折射光。两角均从当地法线量起。
\end{drill}
\fig{fig:d11q}{练习：空气与玻璃界面的两次入射}{
\begin{tikzpicture}
\draw[world](-2,0)--(2,0);\draw[normal](0,-2)--(0,2);
\draw[ray](-1.3,1.3)--(0,0);
\node[above left]at(-1.3,1.3){空气};\node at(-1.4,-1){玻璃};
\node at(-.35,.8){$45^\circ$};\node[below right]at(0,0){$O$};
\begin{scope}[xshift=5.5cm]
\draw[world](-1.5,0)--(1.5,0);\draw[normal](0,-1.8)--(0,1.5);
\draw[ray](-1.149,-.964)--(0,0);\node at(-.35,-.8){$50^\circ$};
\node[above]at(0,1.5){玻璃射向空气};\node at(-1.1,.8){空气};\node at(-1.1,-1.4){$n=\sqrt2$};\end{scope}
\end{tikzpicture}}
\end{question}


\begin{question}
\begin{drill}[平行玻璃板：位移与时间]\label{d:12}
空气中一平行玻璃板厚 \(t=6.00\,\mathrm{cm}\)，折射率 \(\sqrt2\)。光以 \(45^\circ\) 入射。求板内光线的横向位移 \(X\)、出射方向及在玻璃内的传播时间。取 \(c=3.00\times10^8\,\mathrm{m/s}\)。
\end{drill}
\fig{fig:d12q}{练习：平行玻璃板：位移与时间}{
\begin{tikzpicture}
\fill[gray!8](-1,0)rectangle(3,3);
\draw[world](-1,0)--(3,0) (-1,3)--(3,3);
\draw[ray](-1.3,4.3)--(0,3);\draw[normal](0,2)--(0,4);
\node[above right]at(0,3){$P$};\node at(2,1.5){$n=\sqrt2$};
\draw[dim](-.6,0)--(-.6,3)node[midway,left]{$t$};
\end{tikzpicture}}
\end{question}


\begin{question}
\begin{drill}[平面镜上的反射点与总路程]\label{d:13}
同一均匀介质中，平面镜为 \(y=0\)，物点 \(A=(-2,3)\)，接收点 \(B=(4,1)\)，坐标单位均为米。求 \(A\to P\to B\) 的反射点 \(P\) 及总路程，并检查两侧角度。
\end{drill}
\fig{fig:d13q}{练习：平面镜上的反射点与总路程}{
\begin{tikzpicture}
\draw[axis](-3,0)--(5,0)node[right]{$x$};\draw[axis](0,-.3)--(0,3.8)node[above]{$y$};
\fill(-2,3)circle(1.5pt)node[above]{$A(-2,3)$};
\fill(4,1)circle(1.5pt)node[above]{$B(4,1)$};
\foreach\xx in {-2.8,-2.4,...,4.8}{\draw(\xx,0)--++(.18,-.18);}\node[below]at(2,-.3){镜面 $y=0$，反射点 $P$ 待定};
\end{tikzpicture}}
\end{question}

